\chapter{Architecture Overview}

This document studies \textbf{oXiDe SE}, a concrete open source secure-element
operating system and its reference implementation. The architecture presented
here is therefore not only a generic interpretation of GlobalPlatform: it is
the security and software architecture against which the oXiDe SE source tree,
build configurations, and validation results are to be evaluated. Where the
document generalizes from the implementation, it does so explicitly in order to
identify design principles that may be reused by other open secure-element
systems.

oXiDe SE is designed to run under QEMU in order to make development, testing,
and external contributions easier, but its intended deployment targets are
real, readily available microcontrollers. The current ports cover the Arm
Cortex-M architecture profiles Armv6-M, Armv7-M, and Armv8-M. The primary
targets are \texttt{mps2-an385}, \texttt{raspi-pico1}, and
\texttt{raspi-pico2}: the first two can be exercised under QEMU, while the two
Raspberry Pi Pico generations are also validated on widely available physical
boards.

The architecture is also intended to support deployment in the Arm TrustZone
secure world on processors that provide it. In this configuration, oXiDe SE
would expose secure-element services from an isolated execution environment and
act as a \emph{virtual Secure Element} for software running outside that trusted
world.

%--------------------------------------------------
\section{oXiDe SE Architectural Contribution}

GlobalPlatform defines an external management model, its managed entities, and
the protocols used to administer them. It does not prescribe one unique
internal operating-system decomposition. The architectural contribution of
oXiDe SE lies in the way it maps that model onto a small Rust-based embedded
system:

\begin{itemize}
\item a deliberately narrow kernel owns platform-global control flow,
  selection, APDU mediation, isolation, and the structural invariants of the
  persistent registry;
\item Security Domains are first-class administrative application-like objects
  rather than labels or a collection of unrelated special cases embedded
  throughout the kernel;
\item a Security Domain supplies authority, policy, and secure-channel context,
  while the kernel retains validation, non-escalation checks, and final
  mediation of privileged effects;
\item ordinary applications, called \emph{Rustlets}, are independently built
  FAE images executed behind a narrow ABI and a privileged-service boundary.
  They are generally intended for post-issuance deployment, although oXiDe SE
  also permits Rustlets to be embedded in an image before issuance. These two
  positions relative to issuance are shown in
  Figure~\ref{fig:se-lifecycle-actors} and explained further in
  Section~\ref{sec:gp-deployment-workflow};
\item target-specific hardware support and development backends are kept below
  the same APDU and application model, while their distinct security properties
  remain explicit.
\end{itemize}

The central architectural proposition is not merely that GlobalPlatform
commands can be implemented in Rust. It is that GlobalPlatform administration can be decomposed
between a small invariant-enforcing kernel and bounded administrative
applications without giving ordinary applications direct access to
platform-global state. This separation is intended to make authority,
isolation, persistence, and command dispatch independently analyzable.

%--------------------------------------------------
\section{Trusted Computing Base}

For this manual, the trusted computing base (TCB) comprises the components whose
correctness is required to support the security properties sought for an
oXiDe SE image. Its exact hardware portion depends on the selected target, but
its software portion includes:

\begin{itemize}
\item the trusted boot and execution substrate, including the target support
  used for exception handling, memory protection, interrupts, transport,
  persistent flash, cryptographic primitives, and entropy acquisition;
\item the kernel code that parses and mediates APDUs, maintains the selected
  context, enforces isolation, validates privileged requests, owns the registry,
  and commits persistent mutations;
\item every privileged component statically linked into the kernel image,
  including any build-selected kernel-resident module;
\item the root Security Domain and each administrative Security Domain, but only
  for the policy and authority entrusted to that domain. A Rustlet-backed
  Security Domain remains isolated by the kernel, yet its authorization logic
  is part of the policy TCB for its own administrative subtree;
\item the configured trust anchors, privilege assignments, and build-time
  security configuration that determine which authorities the resulting image
  accepts.
\end{itemize}

Ordinary Rustlets are explicitly outside the TCB and must be treated as
potentially malicious. They may affect their own state and obtain only effects
granted through mediated services. A compromised supplementary Security Domain
may misuse its delegated authority, but the kernel is expected to prevent that
authority from escaping its assigned subtree or bypassing platform invariants.

The compiler, image-construction tools, provisioning process, and hardware
supply chain are necessary trust assumptions for producing the initial trusted
image, but they are not resident components of the runtime TCB. QEMU, host-side
test programs, and semihosting services are likewise outside that runtime TCB.
Chapter~2 states the corresponding attacker model, assumptions, and exclusions.

%--------------------------------------------------
\section{How The Architecture Is Assessed In This Manual}

This overview introduces design objectives; it does not present them as
formally established guarantees. The remainder of the manual progressively
describes the mechanisms, intended properties, assumptions, known limits, and
available validation evidence against which oXiDe SE can be assessed:

\begin{itemize}
\item Chapter~2 defines the assets, adversaries, trust assumptions, security
  goals, and explicitly out-of-scope attacks;
\item Chapter~3 defines the APDU, command-dispatch, and secure-channel protocol
  contracts at the external boundary;
\item Chapter~4 defines the managed-object, lifecycle, deployment, and
  administrative-authority invariants inherited from the GlobalPlatform model;
\item Chapter~5 identifies the internal decomposition, kernel responsibilities,
  privileged-service boundary, Security Domain authority, and isolation model;
\item Chapter~6 defines registry persistence and the transactional behavior
  required across interrupted flash updates and resets;
\item Chapter~7 defines the Rustlet execution and ABI model, including the
  application-facing boundary and its isolation requirements;
\item Chapter~8 maps those abstractions to the reference implementation and
  records the present implementation scope and limitations;
\item the appendices refine the T=0 and secure-channel encodings and document
  the provenance and limits of the regression material used for validation.
\end{itemize}

The companion abstract model in \texttt{docs/model.md} collects the principal
state invariants and reasoning obligations in a lightweight Hoare-style
notation. These obligations have not yet been discharged through a complete
formal verification. The source code and executable tests provide implementation
and regression evidence, not a formal proof or a security certification.
Likewise, passing under QEMU does not establish security on physical hardware.

%--------------------------------------------------
\section{Secure Element Ecosystem and Lifecycle}

Before introducing the technical decomposition of the platform, it is useful to
place the secure element in its operational ecosystem. A secure element does
not exist in isolation: it changes hands several times during its life, and
different actors contribute different layers of hardware, software,
configuration and deployed services.

For the purpose of this manual, the surrounding ecosystem is organized around
five actor categories.

\begin{itemize}
\item \textbf{Foundries} manufacture the silicon and deliver the physical secure
  element as a hardware component.
\item \textbf{Providers} integrate the hardware into a usable platform,
  implement the base software stack, and package the secure element together
  with its kernel and foundational services.
\item \textbf{Issuers} distribute the resulting device or object and control its
  initial personalization, provisioning and operational enrollment.
\item \textbf{Service providers} deploy service-specific applications and data,
  either during manufacturing, under agreement with the issuer, or later in the
  field through post-issuance channels such as OTA updates.
\item \textbf{End users} operate the final device and consume the services
  exposed through the secure element, often without directly perceiving the
  existence of the secure element itself.
\end{itemize}

These actors intervene at different moments of the secure-element lifecycle and
do not all manipulate the same interfaces. Foundries first create the physical
substrate. Providers then transform that substrate into a secure-element
platform by activating the hardware, deploying the kernel and installing the
base software that will later anchor the root of trust. Issuers receive that
platform and personalize it for a product line or deployment context. Service
providers add dedicated applications and associated secrets. End users finally
interact with the resulting device through its exposed services and external
communication channels.

\begin{figure}[ht]
  \centering
  \begin{tikzpicture}[
    scale=0.92,
    transform shape,
    stage/.style={
      draw,
      rounded corners=3pt,
      minimum width=2.8cm,
      minimum height=1.3cm,
      align=center,
      font=\small
    },
    state/.style={
      draw,
      dashed,
      rounded corners=3pt,
      minimum width=2.3cm,
      minimum height=0.9cm,
      align=center,
      font=\small,
      fill=white,
      inner sep=3pt
    },
    zone/.style={
      font=\small\itshape
    },
    >=Latex
  ]

    % --- Top line
    \node[stage] (foundry) at (0,0) {\textbf{Foundry}\\Fabrication};
    \node[state] (birth) at (3.3,0) {Birth};
    \node[stage] (provider) at (7.0,0) {\textbf{Provider}\\Activation, kernel\\deployment, base apps};
    \node[state] (rot) at (10.8,0) {Root of Trust};
    \node[stage] (issuer) at (14.6,0) {\textbf{Issuer}\\Personalization, base\\configuration};

    % --- Bottom line (recentered)
    \node[stage] (user) at (7.0,-3.0) {\textbf{End User}\\Service usage, app\\request, field operation};
    \node[state] (prov) at (10.8,-3.0) {Provisioning};
    \node[stage] (sp) at (14.6,-3.0) {\textbf{Service Provider}\\Additional app\\deployment};

    % --- Top flow 
    \draw[->, thick] (foundry.east) -- (birth.west);
    \draw[->, thick] (birth.east) -- (provider.west);
    \draw[->, thick] (provider.east) -- (rot.west);
    \draw[->, thick] (rot.east) -- (issuer.west);

    % --- Separation line (drawn first so nodes can mask it) 
    \draw[dashed, thick] (-1.2,-1.5) -- (16.2,-1.5);

    % --- Zone labels
    \node[zone, anchor=west] at (-1.0,-1.15) {Trusted / controlled environment};
    \node[zone, anchor=west] at (-1.0,-1.85) {Operational / open environment};

    % --- Issuance node (positionné entre les deux)
    \node[state] (issuance) at (10.8,-1.5) {Issuance};

    % --- Arrow 1: Issuer -> Issuance (descend puis gauche)
    \draw[->, thick]
      (issuer.south) -- ++(0,-0.8) |- (issuance.east);

    % --- Arrow 2: Issuance -> End User (gauche puis descend)
    \draw[->, thick]
      (issuance.west) -| (user.north);

    % --- Provisioning flow
    \draw[->, thick] (user.east) -- (prov.west);
    \draw[->, thick] (prov.east) -- (sp.west);

  \end{tikzpicture}
  \caption{Simplified lifecycle of a secure element and transfer of responsibility between actors.}
  \label{fig:se-lifecycle-actors}
\end{figure}

Figure~\ref{fig:se-lifecycle-actors} should not be read as a rigid industrial
workflow. In practice, a single organization may play several roles, and some
deployment steps may occur earlier or later depending on the business model.
The important point for the rest of this manual is that the same secure element
progressively accumulates hardware properties, kernel-level trust anchors,
issuer-specific configuration and service-provider payloads before being used
in the field.

The architecture is structured into four major components.

%--------------------------------------------------
\section{Communication and Protocol}

This component is responsible for handling all interactions between the
external world and the platform. It defines how commands are transported,
interpreted and secured. It is the entry point through which incoming requests
must pass for decoding and validation before execution;
the effectiveness of those checks is assessed in the later protocol and
implementation chapters.

\begin{itemize}
\item \textbf{APDU transport layer} \gpref{gp-card}{§5.3, §9.4}
\item \textbf{APDU parser and dispatcher} \gpref{gp-card}{§9.4, §11}
\item \textbf{Secure Channel Protocol (SCP)} \gpref{scp03}{§4, §6, §7}
\end{itemize}

The APDU transport layer corresponds to the low-level communication defined in
ISO/IEC 7816 and referenced in the GlobalPlatform card specification. Sections
§5.3 and §9.4 describe how commands are encapsulated and exchanged. This layer
is responsible for transporting APDUs across the selected physical or emulated
medium while preserving the command representation expected by the parser.

The APDU parser and dispatcher implements the command interpretation logic
described in §9.4 and §11. It decodes the APDU structure and maps each command
(such as INSTALL, LOAD, DELETE) to its corresponding internal handler. This
component effectively bridges the external protocol with the internal execution
model.

The Secure Channel Protocol, described in SCP03 (§4, §6, §7), introduces
cryptographic protection over APDU exchanges. It defines how sessions are
established, how keys are derived, and how messages are authenticated and
optionally encrypted. The protocol is intended to provide integrity and
confidentiality for administrative commands when it is correctly configured,
implemented, and used within its stated assumptions.

%--------------------------------------------------
\section{GlobalPlatform Conceptual Model}

This component introduces the main conceptual entities manipulated by a
GlobalPlatform secure element. It describes what kinds of objects exist on the
secure element, how they relate to each other, how their lifecycle is controlled,
and how administration authority is represented. It remains at the level
of GlobalPlatform concepts rather than of internal software
decomposition.

\begin{itemize}
\item \textbf{Managed objects}: load files, application instances, and Security
  Domains \gpref{gp-card}{§7, §9}
\item \textbf{AID-based identification and relationships}
  \gpref{gp-card}{§6.2, §7}
\item \textbf{Lifecycle states and transitions}
  \gpref{gp-card}{§9.6}
\item \textbf{Deployment and personalization workflows}
  \gpref{gp-card}{§9, §11}
\item \textbf{Administrative authority and privileges}
  \gpref{gp-card}{§7, §11.1.2}
\end{itemize}

This conceptual layer explains, for instance, why a load file, an
installed application instance, and a Security Domain are distinct
entities even when they are operationally related.

It also explains why commands such as \texttt{LOAD}, \texttt{INSTALL},
\texttt{DELETE}, \texttt{PUT KEY}, or \texttt{SET STATUS} should be read
against a shared administration model rather than as isolated APDU
opcodes.

The implementation-specific mapping of those concepts onto a kernel and
resident applications is deferred to the later execution-model chapter.

%--------------------------------------------------
\section{Model, Security and Core Services}

This component defines the internal structure of the platform. It is
responsible for maintaining consistency, enforcing isolation and providing the
cryptographic foundation required for secure operations.

\begin{itemize}
\item \textbf{Registry of managed objects} \gpref{gp-card}{§9.3}
\item \textbf{AID management} \gpref{gp-card}{§6.2, §9.3.1}
\item \textbf{Secure persistence} \gpref{gp-card}{§9}
\item \textbf{Security Domains} \gpref{gp-card}{§7}
\item \textbf{Privilege and access control} \gpref{gp-card}{§11.1.2}
\item \textbf{Key management} \gpref{scp03}{§6.2, §6.3}
\item \textbf{Cryptographic services} \gpref{scp03}{§6}
\item \textbf{Logging and audit} \gpref{gp-card}{§9}
\end{itemize}

The registry, described in §9.3, maintains the internal representation of all
entities managed by the platform. It is the authoritative source of truth for
application instances, load files and Security Domains.

AID management, defined in §6.2 and §9.3.1, is intended to maintain unique
identification of all entities. It also supports selection mechanisms used by
APDU commands.

Secure persistence is intended to keep the registry and associated metadata
transactionally consistent across resets. The concrete mechanism, failure
model, and current validation evidence are described later in this manual.

Security Domains (§7) define administrative boundaries. They allow grouping
applications under a shared authority and delegating management capabilities.

Privileges and access control (§11.1.2) define what operations an application
or domain can perform. These include administrative rights such as loading or
deleting applications.

Key management and cryptographic services are defined in SCP03 (§6). They
specify how keys are stored, derived and used for secure communication.

Logging and audit mechanisms provide traceability of operations and support
debugging and security analysis.

%--------------------------------------------------
\section{System Decomposition and Execution Model}

This component describes the internal software architecture chosen by the
project to realize the GlobalPlatform concepts introduced above. It is
explicitly an implementation reading of the specification rather than a
normative statement about how every GlobalPlatform secure element must be built.

\begin{itemize}
\item \textbf{A small kernel handling platform-global control flow}
\item \textbf{Administrative applications}, notably Security Domains
\item \textbf{Ordinary applications selected and invoked through APDUs}
\item \textbf{A privileged service boundary} between applications and
  kernel-managed effects
\end{itemize}

The resulting model keeps \texttt{SELECT} and other context-defining
operations in the kernel while routing most other APDUs to the currently
selected AID-associated entity.

In that decomposition, Security Domains are realized as privileged
administrative applications rather than as ad hoc kernel branches for
every management command. This keeps APDU dispatch uniform while still
preserving a strict authority boundary for sensitive operations.

%--------------------------------------------------

\section{Layered Architecture}

The architecture introduced in the previous sections can be summarized as a
layered stack. 

\begin{figure}[ht]
  \centering
  \begin{tikzpicture}[
    layer/.style={
      draw,
      rounded corners=3pt,
      minimum width=10cm,
      minimum height=1.1cm,
      align=center
    },
    >=Latex
  ]
    \node[layer] (apdu) at (0,0)    {\textbf{APDU Communication}\\Transport, APDU parsing, secure channel};
    \node[layer] (admin) at (0,1.8) {\textbf{Administration}\\Loading, installation, lifecycle, deployment};
    \node[layer] (model) at (0,3.6) {\textbf{Model and Security}\\Registry, AID management, domains, privileges, cryptographic services};
    \node[layer] (runtime) at (0,5.4) {\textbf{Runtime and Services}\\Application execution, invocation, service model};
    \node[layer] (apps) at (0,7.2)  {\textbf{Applications}\\Installed and managed application instances};

    \draw[->, thick] (apdu.north) -- (admin.south);
    \draw[->, thick] (admin.north) -- (model.south);
    \draw[->, thick] (model.north) -- (runtime.south);
    \draw[->, thick] (runtime.north) -- (apps.south);
  \end{tikzpicture}
  \caption{Layered view of an open source GlobalPlatform environment architecture.}
  \label{fig:gp-layered-architecture}
\end{figure}

This representation emphasizes the separation of concerns between
communication, administration, internal security mechanisms and application
execution. It also makes explicit the direction of dependency: upper layers
rely on the services provided by lower layers, while lower layers remain
agnostic to the internal logic of upper layers.

Figure~\ref{fig:gp-layered-architecture} highlights the main architectural
principle of the platform. The bottom layer handles communication concerns and
provides a normalized command entry point through APDU processing and secure
messaging. On top of it, the administration layer implements the management
logic required to deploy and control applications. The model and security layer
provides the persistent representation of managed entities together with the
mechanisms that enforce isolation and authorization. Finally, the runtime layer
hosts the execution model exposed to applications. This organization reflects
the fact that GlobalPlatform strongly structures the lower and middle parts of
the stack, while leaving more freedom to the implementation of the execution
environment itself.

%--------------------------------------------------
\section{Glossary}

Unless stated otherwise, the technical vocabulary of this document follows
ISO/IEC~7816 for transport-level concepts and the GlobalPlatform Card
Specification for secure-element-management concepts \gpref{iso7816-4}{§5, §6}
\gpref{gp-card}{§1.4, §6.5, §7}. The purpose of the present glossary is
twofold:
\begin{itemize}
\item to define the technical terms used throughout the document;
\item to keep the text \textbf{secure-element centric}, even though a large
  part of the normative literature remains historically \textbf{card}
  centric.
\end{itemize}

In other words, this document prefers \textbf{secure element} as the generic
system concept, while still acknowledging that ISO/IEC~7816 and
GlobalPlatform often use the word \textbf{card} for the emblematic and
historically dominant form of such a device.

\begin{itemize}
\item \textbf{secure element}: the preferred generic term for the protected
  computing device targeted by this document. It designates the platform that
  stores keys, hosts managed applications, executes APDU commands, and enforces
  the relevant isolation and lifecycle policies.
\item \textbf{card}: the historical and normative term used throughout
  ISO/IEC~7816 and the GlobalPlatform Card Specification for an emblematic
  class of secure element. In this manual, it is retained when the wording of
  the standards is being discussed directly, or when using fixed protocol terms
  such as \emph{card challenge}, \emph{card cryptogram}, or \emph{card
  management command}.
\item \textbf{off-card entity}: the preferred GlobalPlatform term for the
  external party interacting with the secure element.
\item \textbf{interface device}: the ISO/IEC~7816 transport-level term for the
  equipment that exchanges APDUs with the secure element.
\item \textbf{host}: the preferred SCP03 term for the external peer in secure
  channel establishment, notably in \emph{host challenge} and \emph{host
  cryptogram}.
\item \textbf{command APDU}: an APDU sent from the interface device or
  off-card entity to the secure element.
\item \textbf{response APDU}: an APDU returned by the secure element to the
  interface device or off-card entity.
\item \textbf{header}: the four command bytes \texttt{CLA}, \texttt{INS},
  \texttt{P1}, and \texttt{P2}.
\item \textbf{incoming data}: the command payload transported from the
  interface device to the secure element and logically measured by
  \texttt{Lc}.
\item \textbf{outgoing data}: the response payload returned by the secure
  element before the final status word and logically measured by \texttt{Le}.
\item \textbf{status word}: the two-byte trailer \texttt{SW1/SW2} appended to
  a response APDU.
\item \textbf{ATR}: the \emph{Answer To Reset}, i.e. the first byte string
  emitted by the secure element after activation/reset and before the APDU
  command stream begins.
\item \textbf{\texttt{T=0} transport}: the half-duplex, character-oriented
  transport defined by ISO/IEC~7816-3.
\item \textbf{deferred response}: the \texttt{T=0} realization of a logical
  Case~4 exchange, where response data is first announced by \texttt{61xx} and
  then recovered through one or more \texttt{GET RESPONSE} commands.
\item \textbf{application instance}: the preferred term for the installed
  application object managed by the secure element and later selected through
  its AID.
\item \textbf{ordinary application}: a non-administrative application instance
  implementing business logic rather than GlobalPlatform administration.
\item \textbf{administrative application}: an application-level component that
  implements administrative semantics and is granted privileged access to
  kernel-mediated services. In the present architecture, Security Domains are
  realized in this category.
\item \textbf{Security Domain}: the GlobalPlatform administrative principal
  associated with keys, privileges, and management authority.
\item \textbf{Issuer Security Domain}: the distinguished Security Domain that
  anchors the main administrative authority of the secure element.
\item \textbf{Card Manager}: a broader historical label often associated in
  practice with the Issuer Security Domain or with its operational role. This
  document keeps \emph{Issuer Security Domain} as the more precise preferred
  term.
\item \textbf{Executable Load File}: the deployable executable content object
  transported through \texttt{LOAD}. After first introduction, this manual uses
  the shorter form \textbf{load file}.
\item \textbf{GlobalPlatform Registry}: the authoritative catalogue of managed
  objects, identities, privileges, lifecycle states, and dependencies within
  the secure element. After first introduction, this manual uses the shorter
  form \textbf{registry}.
\item \textbf{selected application}: the GlobalPlatform application currently
  selected on a logical channel and therefore designated to receive ordinary
  APDUs after the context-setting phase.
\item \textbf{selected context}: the broader internal execution state that
  includes, but is not limited to, the selected application. This term is kept
  for implementation discussions where the selection state carries additional
  channel or transport information.
\item \textbf{card management command}: the GlobalPlatform term for an APDU
  that operates on managed objects, lifecycle state, key material, or other
  administrative state of the secure element under the authority of a Security
  Domain.
\end{itemize}

%--------------------------------------------------
